\documentclass[10pt,a4paper]{amsart}
\usepackage[margin=19mm]{geometry}
\usepackage{amsmath,amssymb,mathtools}
\usepackage[scr=rsfs,cal=euler]{mathalfa}
\usepackage{xfrac}
\usepackage[unicode,hidelinks,bookmarks=false]{hyperref}
\newcommand{\ie}{i.e.\ }
\newcommand{\cf}{cf.\ }
\newcommand{\resp}{respectively\ }
\newcommand{\Subsection}{\subsection}
\providecommand{\nobreakdash}{\nobreak\hbox{-}}
\setlength{\emergencystretch}{2em}
\multlinegap0pt
\usepackage{float}
\usepackage{bm}
\usepackage{tikz}
\usepackage{mathtools}
\usepackage{enumerate}
\usepackage{xcolor}
\newtheorem{theorem}{Theorem}
\newtheorem{claim}{Claim}
\newcommand{\htheta}{\hat{\theta}}
\title{Fast state transfer via loop weights}
\author{Gabor Lippner, Yujia Shi}
\date{Source: arXiv:2601.20237v1 (CC BY 4.0)}
\begin{document}
\maketitle

\noindent\emph{Source and licence.} Fast state transfer via loop weights. Version arXiv:2601.20237v1 (CC BY 4.0), distributed by arXiv under the Creative Commons Attribution 4.0 International licence, http://creativecommons.org/licenses/by/4.0/, as recorded in the arXiv licence metadata for this exact version and shown on the abstract page https://arxiv.org/abs/2601.20237v1. Full licence notice: \texttt{licenses/arxiv-2601.20237-CC-BY-4.0.txt}.\par\smallskip

\noindent\textit{Editorial scope.} Counted unit: the article's theorem thm:picktheta (source0.tex lines 486--493). The article's complete original proof of it is delivered with it (source0.tex lines 526--560, one \texttt{proof} environment of 35 lines, which the article places away from the statement). No mathematics is written, repaired or re-derived: the statement and the proof are the article's own lines, in the article's own notation, and no retained line is substituted or re-flowed.  Retained uncounted as the source-local context the proof needs: lines 503--507; lines 520--523.  Nothing was omitted from any retained span, so the package carries no omission record.  No retained line contains a citation, so the wrapper carries no bibliography and no bibliography entry.  No retained line contains a figure environment, an image-inclusion command or drawing code, so no image is delivered and the package declares no asset dependency.  Editorial formatting only, in addition: an AMS \texttt{amsart} preamble that restates the article's own declarations and macros used by the retained lines, namely the environments \texttt{theorem}, \texttt{claim} on separate counters, together with the standard packages those lines need.  The retained environments print in this document as Theorem 1, Claim 1 in order of appearance, whereas the article prints the same items with the numbering of its own full document.  The complete native source of the article as distributed in the arXiv e-print for this exact version is bundled unchanged as \texttt{sources/193513/source0.tex}.
\begin{claim}\label{cl:monotone}
    $F(\theta):= \tfrac{-1}{2\cos \theta} + \cos \theta$ is monotone decreasing and convex on $(\pi/2,\pi)$. Furthermore both 
    \[S(\theta)-F(\theta) = -\tan\left(\tfrac{k}{2} \theta\right) \sin \theta\] and 
    \[A(\theta) - F(\theta)= \cot\left(\tfrac{k}{2} \theta\right) \sin \theta\] are monotone decreasing on $(\pi/2,3\pi/4)$ everywhere they are non-singular.
\end{claim}

\begin{claim}\label{cl:F'}
Let $ 0 < c < 1/2$. Then for any $\tfrac{\pi}{2}+\tfrac{c}{\sqrt{k}} \leq x \leq  \tfrac{\pi}{2} + \tfrac{1}{2\sqrt{k}}$ we have 
\[ k \leq |F'(x)| \leq \frac{3k}{c^2}.\]
\end{claim}

\begin{theorem}\label{thm:picktheta} Let $0<c<1/2$ be fixed and let $m$ be an integer such that 
\begin{equation}\label{eq:m_cond} 
k+\frac{2c}{\pi}\sqrt{k}+\frac{1}{2} < 8m+1 < k + \frac{1}{\pi}\sqrt{k}-\frac{1}{2}
\end{equation} 
and denote \(\theta_0 = \tfrac{16m+2}{4k}\pi\) and \(Q = F(\theta_0)\). Then there exists $\theta_1, \theta_2$ such that 
    \begin{align*}\frac{16m+1}{4k} \pi < \theta_1 < \frac{16m+2}{4k}\pi < \theta_2 < \frac{16m+3}{4k}\pi  && \mbox{and} && S(\theta_1)= Q = A(\theta_2).\end{align*}
Furthermore, \[\theta_2-\theta_1 \geq \frac{\sqrt{2}-1}{3k}c^2.\]
\end{theorem}

\begin{proof}[Proof of Theorem~\ref{thm:picktheta}]
Recall that $\theta_0 = \tfrac{8m+1}{2k}\pi$, hence \(\tfrac{k}{2}\theta_0 = 2m\pi + \pi/4\) and thus $\cot(\tfrac{k}{2}\theta_0) = \tan(\tfrac{k}{2}\theta_0)=1$. Then, by the choice of $Q$ we see that 
\[ S(\theta_0 ) = Q -  \sin \theta_0  < Q < Q + \sin(\theta_0) = A(\theta_0).\]
Now let  
\begin{align*} \htheta_1 = \theta_0 - \frac{2}{k}\cdot \frac{\pi}{8} = \frac{16m+1}{4k}\pi && \mbox{and} && \htheta_2 = \theta_0 + \frac{2}{k}\cdot \frac{\pi}{8} = \frac{16m+3}{4k}\pi,\end{align*} and observe that the conditions on $m$ in \eqref{eq:m_cond} imply that 
\begin{align*}\frac{\pi}{2} + \frac{c}{\sqrt{k}} &\leq \htheta_1 & \mbox{and}&& \htheta_2 &\leq \frac{\pi}{2} + \frac{1}{2\sqrt{k}}.
\end{align*} Thus, the conditions of Claim~\ref{cl:F'} are satisfied on the $[\htheta_1,\htheta_2]$ interval. It follows that
\begin{align*}
    F(\htheta_1) &= F(\theta_0)+\int_{\theta_0}^{\htheta_1} F'(x)dx & \mbox{and}& & F(\htheta_2) &= F(\theta_0) + \int_{\theta_0}^{\htheta_2} F'(x)dx \\
    &\geq Q + (\theta_0-\htheta_1)\min_{\htheta_1 \leq x \leq \theta_0} |F'(x)| &&& &\leq Q - (\htheta_2-\theta_0)\min_{\theta_0 \leq x \leq \htheta_2} |F'(x)|\\
    &\geq Q + \frac{\pi}{4k} \cdot k = Q+ \frac{\pi}{4} &&& &\leq Q - \frac{\pi}{4k}\cdot k = Q - \frac{\pi}{4}. \end{align*}
Hence 
\begin{align*}
    S(\htheta_1) &= F(\htheta_1) - \sin \htheta_1 \tan \htheta_1 & \mbox{and}& & A(\htheta_2)&= F(\htheta_2) + \sin \htheta_2 \cot \htheta_2 \\
    &=F(\htheta_1) - \sin \htheta_1 / \sqrt{2}& && &= F(\htheta_2) + \sin \htheta_2/ \sqrt{2} \\
    &\geq F(\htheta_1)-(\sqrt{2}-1) && &&  \leq F(\htheta_2)+\sqrt{2}-1 \\
    &Q + \frac{\pi}{4} - (\sqrt{2}-1) >Q&& && \leq Q - \frac{\pi}{4} + (\sqrt{2}-1) < Q
\end{align*}

Thus, in summary, we see that $S(\htheta_1) > Q > S(\theta_0)$ and $A(\htheta_1) < Q < A(\theta_0)$. Since on the $[\htheta_1, \htheta_2]$ interval both $S$ and $A$ are continuous, there have to exist $\theta_1, \theta_2$ satisfying $\htheta_1 \leq \theta_1 < \theta_0 < \theta_2 \leq \htheta_2$ such that $S(\theta_1) = Q = A(\theta_2)$.

To finish, we need to lower-bound $\theta_2-\theta_1$. We can do that by showing that $F(\theta_1) - F(\theta_2)$ cannot be too small. Using the monotonicity from Claim~\ref{cl:monotone} we can estimate
\begin{align*}F(\theta_1) &= S(\theta_1) + \sin \theta_1 \tan\left(\tfrac{k}{2}\theta_1 \right) 
&\mbox{and}&& F(\theta_2) &= A(\theta_2) - \sin \theta_2 \cot\left(\tfrac{k}{2}\theta_2 \right)
\\ 
&= Q + \sin \theta_1 \tan\left(\tfrac{k}{2}\theta_1 \right) & & &  &= Q - \sin \theta_2 \cot\left(\tfrac{k}{2}\theta_2 \right)
\\ 
&\geq Q + \sin \htheta_1 \tan\left(\tfrac{k}{2}\htheta_1 \right) & & & &\leq Q - \sin \htheta_2 \cot\left(\tfrac{k}{2}\htheta_2 \right)
\\ 
&\geq Q + (\sqrt{2}-1)/2 
&&&   &\leq Q - (\sqrt{2}-1)/2.\end{align*}

So we get, using the upper bound from by Claim~\ref{cl:F'}, that
\[ \sqrt{2}-1 \leq F(\theta_1) - F(\theta_2) \leq (\theta_2 - \theta_1)\max\{|F'(x)|: \theta_1 \leq x \leq \theta_2\} \leq (\theta_2 -\theta_1) \frac{3k}{c^2},\] from which the claimed bound on $\theta_2 - \theta_1$ follows immediately.
\end{proof}

\end{document}
